% Chapter 8: API reference (practice author). PLAN.md section 1.4. Every statement is checked against
% include/numerixx/deriv.hpp and include/numerixx/math.hpp. Tables: T2 tab:api-names, T3 tab:api-methods.
% Example: E16 (code/08-api-usage.cpp); E11's names snippet.
\chapter{API reference}\label{ch:api}

This chapter lists every public name of \code{nxx::deriv}, and those of the shared header
\code{<numerixx/math.hpp>} that it uses, with its declaration, parameters, preconditions, result, failures and cost,
and a pointer to a compiled example. It describes the headers as they are, in their own wording where it summarises a
comment of the header. Part~I explains the methods, \cref{ch:using} shows them in use, and \cref{ch:tables,ch:guarantees} give the
numbers.

\section{Synopsis}\label{sec:api-synopsis}

The module is one header, \code{<numerixx/deriv.hpp>}, which includes \code{<numerixx/math.hpp>}. Both are
header-only, need C++23 and no other library; the CMake target is \code{numerixx::numerixx}. \Cref{tab:api-names}
lists the public names. Everything in a namespace \code{detail} is internal: it may change without notice, and the
book reads it only to generate the tables of \cref{ch:tables}. The enumeration \code{error} may gain values in a
later version. A program usually starts with
\snippet{07-using-expected}{names}

\begin{table}
\sidecaption{The public names. \code{deriv.hpp} declares those in \code{nxx::deriv}, \code{math.hpp} those in
\code{nxx} and \code{nxx::math}.\label{tab:api-names}}
\centering\footnotesize
\begin{tabular}{@{}>{\raggedright}p{45mm}>{\raggedright}p{47mm}l@{}}
\toprule
Name & Kind & Section\\
\midrule
\code{derivative} & function template: two overloads, each with a deprecated twin & \ref{sec:api-derivative}\\
\code{Estimate<T>} & class template, an aggregate & \ref{sec:api-estimate}\\
\code{error} & scoped enumeration: \code{non_finite} & \ref{sec:api-error-order}\\
\code{order} & scoped enumeration over \code{int}: \code{first} to \code{fourth} & \ref{sec:api-error-order}\\
\code{Central<P>}, \code{Forward<P>}, \code{Backward<P>} & class templates, empty method tags &
  \ref{sec:api-methods}\\
\code{Ridders} & class, an empty method tag & \ref{sec:api-methods}\\
\code{NUMERIXX_NO_PRECISION_WARNING} & macro, defined by the user & \ref{sec:api-warning}\\
\code{nxx::real} & concept & \ref{sec:api-real}\\
\code{nxx::math::abs}, \code{isfinite} & function templates & \ref{sec:api-real}\\
\code{nxx::math::root_epsilon} & function template & \ref{sec:api-real}\\
\bottomrule
\end{tabular}
\end{table}

\section{\texorpdfstring{\code{derivative}}{derivative}}\label{sec:api-derivative}

\paragraph{Declarations.} Two overloads, with a step and with the default step:
\begin{signature}
\code{template<typename M = Central<2>, auto N = 1, real T, typename F>}\\
\sigindent\code{requires detail::method<M> && detail::valid_order<N>}\\
\sigindent\sigindent\code{&& detail::callable<F, T>}\\
\code{[[nodiscard]] constexpr std::expected<Estimate<T>, error>}\\
\sigindent\code{derivative(F &&fn, T const &x, std::type_identity_t<T> const &h);}\\[1ex]
\code{template<typename M = Central<2>, auto N = 1, real T, typename F>}\\
\sigindent\code{requires detail::method<M> && detail::valid_order<N>}\\
\sigindent\sigindent\code{&& detail::callable<F, T>}\\
\code{[[nodiscard]] constexpr std::expected<Estimate<T>, error>}\\
\sigindent\code{derivative(F &&fn, T const &x);}
\end{signature}
Each has a twin with the additional constraint \code{detail::noisy<N, T>}, marked \code{[[deprecated]]}
(\cref{sec:api-warning}); it returns the same result.

\paragraph{Template parameters.}
\begingroup\raggedright
\begin{itemize}
  \item \code{M}, the method: \code{Central<P>} (the default is \code{Central<2>}), \code{Forward<P>},
    \code{Backward<P>} or \code{Ridders}. \code{detail::method<M>} accepts the three stencil families with an accepted
    order and \code{Ridders}, and nothing else: the set is closed, so users cannot add a method.
  \item \code{N}, the derivative order: 1 (the default) to 4, as a value of any integer type except \code{bool}, or as
    an \code{order}. \code{detail::valid_order<N>} compares an integer in its own type, so a value beyond the range of
    \code{int} is not wrapped into 1 to 4.
  \item \code{T}, the real type, deduced from \code{x} alone; it must satisfy \code{nxx::real}
    (\cref{sec:api-real}). The type of the step is non-deduced, \code{std::type_identity_t<T>}, so that a
    \code{double} step converts to a multiprecision \code{T} instead of making deduction fail.
  \item \code{F}, the callable. The constraint \code{detail::callable<F, T>} is
    \code{std::invocable<F &, T const &>}, with a result convertible to \code{T}: \code{fn} is called as a named
    lvalue through \code{std::invoke}, with a \code{T const} lvalue, so member pointers work, and so does a callable
    that rejects rvalues.
\end{itemize}
\endgroup

\paragraph{Parameters.} \code{fn}, the function to differentiate. \code{x}, the point. \code{h}, the step of a
stencil, or the initial step of Ridders' method.

\paragraph{Preconditions.} \code{h > 0} and \code{h} finite, checked with \code{assert}: a bad step is a bug in the
caller, not a run-time failure. For a stencil, \code{h} well above $\ulp(x)$; a step below the resolution of
\code{x} is detected and returns \code{non_finite}, but a step of a few ulps gives a finite value that is mostly
rounding. The overload without a step has no precondition.

\paragraph{Default steps.} For a stencil of accuracy order P, $h = \epsmach^{1/(P+N)} \max(\abs{x}, 1)$, computed as
\code{root_epsilon<P + N, T>()} times the scale. For Ridders' method the initial step is $\max(\abs{x}, 1)/32$ for
$f'$ and $\max(\abs{x}, 1)/4$ for $f''$ to $f''''$. The header advises passing a step when \code{fn} varies faster
than $\max(\abs{x}, 1)$, or when the derivative is much smaller than the error $\epsmach^{P/(P+N)}$, and using the
overload without a step in generic code that varies N (\cref{tab:default-steps}).

\paragraph{Result.} An \code{Estimate<T>} (\cref{sec:api-estimate}) whose \code{abs_error} is filled by Ridders'
method and empty for the stencils, or \code{error::non_finite}.

\paragraph{Failures.} \code{error::non_finite} when no finite estimate exists: \code{x} is not finite; \code{fn}
returned an infinite or NaN value; a difference, term, sum or quotient would overflow; a point $x + k h$ would
overflow; or \code{h} is below the resolution of \code{x} (for a rule with more than two nodes, also when neighbouring
points round to the same value). For Ridders' method, specifically: \code{fn(x)} is not finite for an even N (after
one evaluation), or no tableau reached two finite levels. A level at which \code{fn} is not finite makes the method
restart below that step; if the levels run out during restarts after an untrusted result, that result is returned, as
a value with its large \code{abs_error}, not as a failure.

\paragraph{Cost.} A stencil calls \code{fn} once per node: $P + N - 1$ times for \code{Central<P>}, $P + N$ for
\code{Forward<P>} and \code{Backward<P>}. Ridders' method calls it 2 times per level for $f'$ and $f''$, 4 times for
$f'''$ and $f''''$, plus \code{fn(x)} once for $f''$ and $f''''$: at most 64, 65, 128 and 129 times
(\cref{tab:cost}). \code{fn} is called only at finite points. No memory is allocated.

\paragraph{Exceptions and constant evaluation.} Numerixx throws nothing and catches nothing; an exception thrown by
\code{fn} or by the operations of \code{T} passes through unchanged. Both overloads are \code{constexpr}. Failures
are detected before the arithmetic that would produce them, so they are error values in constant expressions too,
except an overflow of a final quotient, which GCC rejects at compile time (\cref{sb:constexpr-overflow}).

\paragraph{Example.} Every accepted form, on the thread's enthalpy at $T_0$ (\code{code/08-api-usage.cpp}):
\snippet{08-api-usage}{api-overloads}

\section{\texorpdfstring{\code{Estimate<T>}}{Estimate<T>}}\label{sec:api-estimate}

\begin{signature}
\code{template<real T> struct Estimate}~\texttt{\{}\\
\sigindent\code{T value;}\\
\sigindent\code{std::optional<T> abs_error;}\\
\sigindent\code{constexpr operator T() const;}\\
\sigindent\code{friend constexpr bool}\\
\sigindent\sigindent\code{operator==(Estimate const &, Estimate const &) = default;}\\
\texttt{\};}
\end{signature}
\code{value} is the estimate of the derivative. \code{abs_error} is an estimate of $\abs{\text{value} -
\text{exact}}$, not a guaranteed bound; it is empty for the stencils. For Ridders' method it is the change of the best
tableau entry from its parent, but at least the rounding error of the difference it came from and at least its distance
to the latest diagonal entry (\cref{sb:floor-coverage}).

\code{Estimate} is an aggregate, so it takes designated initialisers and structured bindings,
\code{auto const [value, abs_error] = *r}. It converts implicitly to \code{T}. Where a conversion is not considered,
in \code{auto}, \code{std::format} and the arithmetic operator templates of Boost.Multiprecision (its comparisons do
accept it), write \code{e.value} or \code{T(e)}. The conversion also applies where it may surprise: \code{if (*r)} tests the value against 0, not whether
there is a result, and \code{int n = *r} truncates the value (\cref{sec:using-result}). The defaulted \code{==}
compares both members; without it, \code{==} would compare only the values, through the conversion.

\section{\texorpdfstring{\code{error}}{error} and \texorpdfstring{\code{order}}{order}}\label{sec:api-error-order}

\begin{signature}
\code{enum class error { non_finite };}\\
\code{enum class order : int { first = 1, second = 2, third = 3, fourth = 4 };}
\end{signature}
\code{error} says why \code{derivative} returned no estimate; its one value is described in
\cref{sec:api-derivative}. More values may be added, so a \code{switch} on it should have a \code{default}.
\code{order} names the derivative orders for readability: \code{derivative<Ridders, order::second>} is
\code{derivative<Ridders, 2>}. It is scoped, so its names stay out of \code{nxx::deriv} and it does not convert to
\code{int} elsewhere.

\section{The methods}\label{sec:api-methods}

\begin{signature}
\code{template<int P> struct Central;}\hfill\code{// even P, 2 <= P <= 10}\\
\code{template<int P> struct Forward;}\hfill\code{// 1 <= P <= 6}\\
\code{template<int P> struct Backward;}\hfill\code{// 1 <= P <= 6}\\
\code{struct Ridders {};}
\end{signature}
The methods are empty tags, passed as the first template argument of \code{derivative}; their weights are internal.
P is the accuracy order, the $\order(h^P)$ of the error, not the number of points. A tag with an order outside its
range is rejected by the constraint of \code{derivative}; completing the type itself fails a \code{static_assert}
(``\code{Central<P> needs an even order P from 2 to 10}'', ``\code{Forward<P> needs an order P from 1 to 6}''). A
stencil whose integer weights the type \code{T} cannot hold exactly fails a \code{static_assert} too (``this rule's
integer weights are not exact in T''); every type with 24 significant bits or more holds them all.
\Cref{tab:api-methods} gives the nodes, the cost and the default step of each.

\begin{table}
\sidecaption{The methods for the N-th derivative: nodes in units of $h$, evaluations of \code{fn}, default step
and error estimate, as \code{deriv.hpp} documents them. $s = \max(\abs{x}, 1)$. The weights are in
\cref{ch:tables}.\label{tab:api-methods}}
\centering\footnotesize
\begin{tabular}{@{}>{\raggedright}p{17mm}>{\raggedright}p{30mm}>{\raggedright}p{19mm}>{\raggedright}p{14mm}%
>{\raggedright\arraybackslash}p{12mm}@{}}
\toprule
Method & Nodes & Evaluations & Default step & Estimate\\
\midrule
\code{Central<P>} & the $P + N - 1$ integers around 0; 0 included for even N only & $P + N - 1$ &
  $\epsmach^{1/(P+N)} s$ & no\\
\code{Forward<P>} & $0, 1, \ldots, P + N - 1$ & $P + N$ & $\epsmach^{1/(P+N)} s$ & no\\
\code{Backward<P>} & $-(P + N - 1), \ldots, 0$ & $P + N$ & $\epsmach^{1/(P+N)} s$ & no\\
\code{Ridders} & \code{Central<2>} of order N at $h/2^i$: $\pm h_i$ for $N \le 2$, $\pm h_i, \pm 2 h_i$ for $N \ge
  3$, and $x$ for even N & 2 per level ($N \le 2$), 4 ($N \ge 3$), plus $f(x)$ once for even N; at most 64, 65, 128,
  129 & $s/32$ ($N = 1$), $s/4$ ($N \ge 2$) & yes\\
\bottomrule
\end{tabular}
\end{table}

The two-node rules, \code{Central<2>} for $f'$, \code{Forward<1>} and \code{Backward<1>}, divide by the distance
between the points actually evaluated, and \code{Central<2>} for $f''$ uses the actual offsets of its points
(\cref{sb:actual-distance,sb:actual-offsets}); rules with more nodes divide the weighted sum by $\den h$ and then by
$h$, $N - 1$ more times, never by $h^N$. \code{Ridders} extrapolates the \code{Central<2>} rule of order N in
a Richardson tableau until rounding stops the improvement; it restarts below a step at which \code{fn} is not finite,
or below one whose result has at most two correct bits by its own estimate while that estimate is far above rounding.
Its constants are in \cref{tab:ridders-constants}.

\section{The precision warning}\label{sec:api-warning}

A fourth derivative in a type whose machine epsilon is above $10^{-7}$ is mostly rounding noise. The test is made with
integers, $\beta^{\text{digits} - 1} < 10^7$ with $\beta$ the radix: \code{float}, \code{decimal32} and
\code{bfloat16} qualify, \code{double} does not. Such a call selects the deprecated twin of \code{derivative}, whose
warning reads ``the 4th derivative in a type with 7 or fewer significant digits is mostly rounding noise: use double,
or silence this warning at the call or with NUMERIXX\_NO\_PRECISION\_WARNING''. The warning is reported at the
caller's line, so a pragma there silences it (\code{-Wdeprecated-declarations} on GCC and Clang, C4996 on MSVC), and a
system-header include does not hide it. Defining \code{NUMERIXX_NO_PRECISION_WARNING} turns it off everywhere; it must
be defined alike in every translation unit, best on the command line, since overload resolution would otherwise differ
between them. \Cref{sec:orders-float} shows the pragma in a compiled example.

\section{\texorpdfstring{\code{nxx::real}}{nxx::real} and \texorpdfstring{\code{nxx::math}}{nxx::math}}\label{sec:api-real}

\begin{signature}
\code{template<typename T> concept real = /* see below */;}\\[1ex]
\code{namespace nxx::math}~\texttt{\{}\\
\sigindent\code{template<real T> constexpr T abs(T const &x);}\\
\sigindent\code{template<real T> constexpr bool isfinite(T const &x);}\\
\sigindent\code{template<int N, real T> constexpr T root_epsilon();}\\
\texttt{\}}
\end{signature}
\code{real<T>} holds when \code{std::numeric_limits<T>} is specialised with \code{is_integer} false,
\code{is_bounded} true and \code{is_exact} false, so that \code{max()} is finite and \code{epsilon()} positive; when
\code{T} is constructible from \code{int} and totally ordered; and when $a + b$, $a - b$, $a \times b$, $a / b$ and
$-a$ convert to \code{T}, which admits expression templates (\cref{sb:expression-templates}). The tests check that
\code{float}, \code{double}, \code{long double}, \code{cpp_bin_float_50} with and without expression templates and
\code{cpp_dec_float_50} satisfy it, and that \code{int}, \code{unsigned} and \code{cpp_rational} do not.

The helpers use comparisons and $+ - \times /$ only, so they run in constant expressions and accept user types.
\code{abs(x)} returns $\abs{x}$. \code{isfinite(x)} is false for infinities and NaN: it tests
\code{x == x && abs(x) <= max()}. \code{root_epsilon<N, T>()} returns $\epsmach^{1/N}$ for $N \ge 2$ (a
\code{static_assert}), accurate to within an ulp, by Newton's iteration on epsilon scaled into $(\beta^{-N}, 1]$; for
\code{double}, \code{root_epsilon<3, double>()} is \num{\valRootEpsThreeDouble}, the default step of
\code{Central<2>} for $f'$ at $\abs{x} \le 1$.

\section{Rejected calls}\label{sec:api-rejected}

What the constraints reject does not compile: the error is reported at the call, not at run time. The example tests
each case with a \code{requires}-expression:
\snippet{08-api-usage}{api-rejected}
A type given explicitly as the second template argument, as in
\code{derivative<Ridders, double>(fn, x)}, does not compile either, because that argument is the order; \code{T} is always deduced from \code{x}. A callable that cannot
be called with a \code{T const} lvalue, or whose result does not convert to \code{T}, fails \code{detail::callable}.
